% CONSERVATIVE REVIEW DRAFT -- replaces the EM-interaction subsection only.
% Replace from \subsection{Interaction of Matter with Electromagnetic Radiation}
% through its table, stopping before the VELOCITY GAUGE placeholder.
% The parent theory introduction and the velocity-gauge subsection are not included.
% New/replacement wording and corrected mathematical terms use \myblue{...}.
% Retained prose/equations remain black; typographical cleanup is not marked.
% Requires the existing \myblue and \vb macros, and amsmath.
%
% REVIEW NOTES
% - Original paragraph sequence retained: TDKS -> plane wave -> fields ->
%   Coulomb gauge -> gauge freedom -> length gauge -> kick -> observables ->
%   finite-q probes -> applications/table.
% - Five-column implementation-oriented table includes a fifth magnetic/spin-probe row.
% - Scalar-potential probes remain electromagnetic and have gauge freedom.
% - q_e=-e in Gaussian units; atomic units begin in the implementation.
% - Removed a repeated finite-system dipole sentence and the incorrect
%   IXS/magnon assignments in the transverse-vector-potential table row.
% - Detailed velocity-gauge transformation and RMG kick normalization are deferred.
% - Updated application/reference placements: ARPES and demagnetization in VG;
%   nonresonant IXS in the scalar row; magnons in a separate paragraph.
% - NEW citation keys below: copy these entries into the master .bib.
%   Existing magnon/demagnetization keys are retained from the original source.
%
% @article{IsmailBeigi-2001-EM,
%   author  = {Ismail-Beigi, Sohrab and Chang, Eric K. and Louie, Steven G.},
%   title   = {Coupling of Nonlocal Potentials to Electromagnetic Fields},
%   journal = {Physical Review Letters},
%   year    = {2001},
%   volume  = {87},
%   pages   = {087402},
%   doi     = {10.1103/PhysRevLett.87.087402}
% }
%
% @article{Bonafe-Rubio-2025,
%   author = {Bonaf{\'e}, Franco P. and Albar, Esra Ilke and Ohlmann, Sebastian T. and Kosheleva, Valeriia P. and Bustamante, Carlos M. and Troisi, Francesco and Rubio, Angel and Appel, Heiko},
%   title = {Full minimal coupling {Maxwell-TDDFT}: An ab initio framework for light-matter interaction beyond the dipole approximation},
%   journal = {Physical Review B}, year = {2025}, volume = {111},
%   pages = {085114}, doi = {10.1103/PhysRevB.111.085114}
% }
% @article{Rubio-angular-2017,
%   author = {De Giovannini, Umberto and H{\"u}bener, Hannes and Rubio, Angel},
%   title = {A First-Principles Time-Dependent Density Functional Theory Framework for Spin and Time-Resolved Angular-Resolved Photoelectron Spectroscopy in Periodic Systems},
%   journal = {Journal of Chemical Theory and Computation}, year = {2017},
%   doi = {10.1021/acs.jctc.6b00897}
% }
% @article{Dixit-2014-IXS,
%   author = {Dixit, Gopal and Slowik, Jan Malte and Santra, Robin},
%   title = {Theory of time-resolved nonresonant x-ray scattering for imaging ultrafast coherent electron motion},
%   journal = {Physical Review A}, year = {2014}, volume = {89},
%   pages = {043409}, doi = {10.1103/PhysRevA.89.043409}
% }
%
% MERGE NOTES
% - Ensure the following subsection defines \label{sec:velocity_gauge}.
% - The theory-level introduction should agree with the updated notation.tex:
%   Gaussian units with explicit constants; implementation in atomic units.
% - The B column distinguishes electromagnetic fields from the auxiliary spin field.
% - If this table is placed in a two-column layout, use table*.

%\subsection{Interaction of Matter with Electromagnetic Radiation}
\subsection{External Perturbations in RT-TDDFT: Gauges, Probes, and Spectroscopies}
\label{sec:em-interaction}

The starting point for the present formulation is the time-dependent
Kohn--Sham equation for single-particle orbitals $\Psi_j(\vb r,t)$ in the
presence of external electromagnetic fields which, within the
minimal coupling framework, takes the form
\begin{equation}
 \imath\hbar\frac{\partial}{\partial t}\Psi_j(\vb r,t)
 =\left[
 \frac{1}{2m_e}
 \left(\vb{\hat p}-\frac{q_e}{c}\vb A_{\mathrm{ext}}(\vb r,t)\right)^2
 +q_e\Phi_{\mathrm{ext}}(\vb r,t)
 +\hat V_{\mathrm{KS}}[\rho](t)
 \right]\Psi_j(\vb r,t).
 \label{eq:em-tdks}
\end{equation}
Here $\vb A_{\mathrm{ext}}$ and $\Phi_{\mathrm{ext}}$ are vector and scalar
potentials representing the external electromagnetic (EM) field,
$\vb{\hat p}=-\imath\hbar\nabla$ is the canonical momentum operator,
 $q_e=-e$, with $e>0$, is the electron charge, and $\rho$ denotes electron number density.
We use Gaussian electromagnetic units, retaining $\hbar$, $m_e$, $e$,
and $c$ explicitly. 
The Kohn--Sham potential, $\hat V_{\mathrm{KS}}[\rho](t)$, contains the
ionic potential, Hartree potential, and exchange-correlation potential:
\begin{equation}
 \hat V_{\mathrm{KS}}
 =\hat V_{\mathrm{ion}}+\hat V_{\mathrm H}[\rho](\vb r,t)
 +\hat V_{\mathrm{xc}}[\rho](\vb r,t).
\end{equation}
In the present implementation based on norm-conserving
pseudopotentials, $\hat V_{\mathrm{ion}}$ contains local and
nonlocal parts. Gauge transformation of the Hamiltonian may introduce additional  terms 
that  must be accounted  for in both, Hamiltonian and  the current density 
operator.\cite{pickard-mauri-2001,Mattiat-Luber-2022,IsmailBeigi-2001-EM}
The working model uses a scalar adiabatic exchange-correlation
potential. Explicit exchange-correlation vector potentials and spin coupling
to an applied magnetic field are not included. 
The choice of the specific form of $\vb A_{\mathrm{ext}}$ and
$\Phi_{\mathrm{ext}}$ depends on the problem and is subject to gauge
freedom.\cite{gauge-theory}


For a monochromatic plane-wave component of the external EM field,
the vector potential may be represented as
\begin{equation}
 \vb A_{\mathrm{ext}}(\vb r,t)
 =\vb A_0 e^{\imath(\vb q\cdot\vb r-\omega t)}+\mathrm{c.c.},
 \label{eq:plane_wave_vector_potential}
\end{equation}
where c.c. denotes complex conjugation ensuring that $\vb A_{\mathrm{ext}}(\vb r,t)$ is real valued,
$\omega$ is the angular frequency, the wave vector $\vb q$
describes the propagation direction with the corresponding  wavelength
$\lambda=2\pi/|\vb q|$, and $\vb A_0=[A_x^0,A_y^0,A_z^0]$
is a complex vector  amplitude.
The electric and magnetic fields are
expressed in terms of the vector and scalar potentials as
\begin{align}
 \vb E(\vb r,t)
 &=-\nabla\Phi_{\mathrm{ext}}(\vb r,t)
   -\frac{1}{c}\frac{\partial\vb A_{\mathrm{ext}}(\vb r,t)}{\partial t},
 %\label{eq:em_fields_from_potentials}\\
 \label{eq:em_E_from_A+Phi}\\
 \vb B(\vb r,t)&=\nabla\times\vb A_{\mathrm{ext}}(\vb r,t).
 \label{eq:em_B_from_A}
\end{align}

For each nonzero Fourier wave vector $\vb q$, which specifies
the direction of spatial variation of the perturbation,
a vector-field amplitude $\vb A_0$ can be decomposed into a \emph{transverse}
(divergence-free) component perpendicular to $\vb q$ and a
\emph{longitudinal} (curl-free) component parallel to $\vb q$.

Choosing  the Coulomb gauge,
\begin{equation}
  \nabla\cdot\vb A_{\mathrm{ext}}(\vb r,t)=0,
  \label{eq:coulomb_gauge_condition}
\end{equation}
ensures that, 
for each nonzero Fourier wave vector $\vb q$,
the vector potential $\vb A_{\mathrm{ext}}(\vb r,t)$ 
is purely transverse  and  carries only the transverse
component of the electric field $\vb E_{\perp}(\vb r,t)$
through the second (time-derivative) term,
$-(1/c)\partial_t\vb A_{\mathrm{ext}}$,
in Eq.~\eqref{eq:em_E_from_A+Phi}, as well as it carries  the magnetic field
$\vb B(\vb r,t)$ through
$\nabla\times\vb A_{\mathrm{ext}}$
(Eq.~\eqref{eq:em_B_from_A}).
% 
For a plane wave, this becomes apparent upon applying the Coulomb
gauge Eq.~\eqref{eq:coulomb_gauge_condition} to
Eq.~\eqref{eq:plane_wave_vector_potential}, which yields
the transversality condition
\begin{equation}
 \vb q\cdot\vb A_0=0,
 \label{eq:transversality_condition}
\end{equation}
requiring that the polarization of the radiation field is perpendicular
to its propagation direction $\vb q$, for all $\vb q\neq 0$.
%
%Similarly, for $\vb q\neq 0$ the contribution to electric field, $\vb E(\vb r, t)$, due
%to the scalar potential (1st term,  $-\nabla\Phi_{\mathrm{ext}}$, in Eq.~\eqref{eq:em_E_from_A+Phi})
%is  exclusively longitudinal in nature  $\vb E(\vb r, t)= \vb E_{\parallel}(\vb r,t)$. 
%This can be seen by noticing that the  Fourier transform vector of $\nabla\Phi_{\mathrm{ext}}$ with respect to $\vb r$
%is parallel to $\vb q$ for every  $\vb q$:  
%$ \mathcal F [ \nabla\Phi_{\mathrm{ext}}] = \imath  \vb q \tilde  \Phi_{\mathrm{ext}}(\vb q, t)$, where the Fourier transform relation 
%$ \mathcal F [ \nabla  f(\vb r) ]  \to  \imath  \vb q  \tilde f(\vb q) $ was used.
%Also, the transverse projection $(1 -\hat {\vb  q} \hat {\vb  q} ) \cdot \imath \vb q \tilde \Phi_{\mathrm{ext}}$ vanishes identically.
%

Similarly, for $\vb q\ne0$ the contribution to
the electric field $\vb E(\vb r,t)$ due to the scalar
potential (first term, $-\nabla\Phi_{\mathrm{ext}}$,
in Eq.~\eqref{eq:em_E_from_A+Phi}) is exclusively longitudinal
in nature
%\myblue{; in the Coulomb gauge, this contribution equals $\vb E_{\parallel}(\vb r,t)$}.
This can be seen by noticing that the Fourier transform
of $-\nabla\Phi_{\mathrm{ext}}$ with respect to $\vb r$
is parallel to $\vb q$ for every nonzero $\vb q$:
\[
  \mathcal F_{\vb r}[-\nabla\Phi_{\mathrm{ext}}]
  =-i\vb q\,\widetilde{\Phi}_{\mathrm{ext}}(\vb q,t),
\]
where the Fourier transform relation
$\mathcal F[\nabla f(\vb r)]=i\vb q\,\widetilde f(\vb q)$
was used.
Also, the transverse projection
$(\mathbf I-\hat{\vb q}\hat{\vb q})\cdot
[-i\vb q\,\widetilde{\Phi}_{\mathrm{ext}}]$
vanishes identically, where
$\hat{\vb q}=\vb q/|\vb q|$ and $\mathbf I$ is the identity tensor.


That is, for $\vb q\neq 0$, the scalar potential  $\Phi_{\mathrm{ext}}$ contributes only  to $\vb E_\parallel$ in any gauge, 
while $\vb A_{\mathrm{ext}}$ contributes  only to $\vb E_\perp$  as a result of the Coulomb gauge.
For the spatially uniform component ($\vb q=0$), $\hat{\vb q}$ is undefined, and
there is no wave-vector direction relative to which a field
can be classified as longitudinal or transverse;
accordingly, Eq.~\eqref{eq:coulomb_gauge_condition} places no constraint on the uniform
component of $\vb A_{\mathrm{ext}}$, since
$\imath\vb q\cdot\tilde{\vb A}_{\mathrm{ext}}=0$ holds trivially at $\vb q=0$.
%This residual freedom  for $\vb q= 0$  within the Coulomb gauge relates the length and
%velocity gauges used in RT-TDDFT, introduced below.}
For $\vb q= 0$, the residual freedom within the Coulomb gauge allows for the  freedom between \emph{length} and \emph{velocity}
gauges used rt-TDDFT.


In general, the electromagnetic potentials ($\vb A$ and $\Phi$) and
the KS potential all appear simultaneously. However, gauge freedom
allows the electromagnetic potentials to be transformed to a
representation suitable for a given problem, using
an arbitrarily chosen  scalar function $\chi(\vb r,t)$
and the gauge transformation rules
\begin{align}
 \vb A&\to\vb A+\nabla\chi(\vb r,t),\\
 \Phi&\to\Phi-\frac{1}{c}
 \frac{\partial\chi(\vb r,t)}{\partial t},
 \label{eq:gauge-rules}
\end{align}
which  guarantee that the fields $\vb E$ and $\vb B$ are unchanged.
A specific choice of $\chi$ relevant to periodic systems, and corresponding orbital and hamiltonian transfromation  are  discussed
in Sec.~\ref{sec:velocity_gauge}.

In finite-system optical applications, a common choice for
RT-TDDFT is the \emph{length gauge} within the electric-dipole
approximation. This approximation applies in the long-wavelength
limit, $|\vb q|L\ll1$, where the radiation wavelength
$\lambda=2\pi/|\vb q|$ is much larger than the characteristic
size $L$ of the system. The applied electric field $\vb E(t)$
is then treated as spatially uniform over the system and represented  by $\vb A=0$ 
 and the dipolar scalar potential
\begin{equation}
  \Phi(\vb r,t)=-\vb E(t)\cdot\vb r,
  \label{eq:EM_dipolar}
\end{equation}
as can be  verified by  application of Eqs.~\eqref{eq:em_E_from_A+Phi}  and   \eqref{eq:em_B_from_A}. %  from $\vb E=-\nabla\Phi$.
For periodic systems, however, the scalar-potential form in  Eq.~\eqref{eq:EM_dipolar}
 is not compatible with periodic boundary conditions, resulting in 
 a sawtooth potential with a discontinuity at the cell boundary. This motivates  
the use of  \emph{velocity-gauge} formulation, in which
the $ \Phi(\vb r,t)=0$  and the  external field $\vb E(t)$ enters through a spatially uniform  vector potential
$\vb A(t)$ via  time-derivative term of Eq.~\eqref{eq:em_E_from_A+Phi}.

While the plane-wave form Eq.~\eqref{eq:plane_wave_vector_potential}
describes a monochromatic perturbation, in practice
the homogeneous electric field can be assigned a general
temporal profile.
A particularly useful choice is the Dirac delta kick, 
 $\vb E(t)=\hat{\vb e} \mathcal E_0\delta(t)$, where
$\hat{\vb e}$ is the polarization direction and $\mathcal E_0$ is  a scalar  perturbation strength.
Since the delta function has a flat spectral content, 
it delivers a broadband impulsive perturbation of magnitude $\mathcal E_0$ at $t=0$, 
The frequency-dependent absorption spectrum is then obtained from
the Fourier transform of the induced time-dependent observable
monitored during the propagation (i.e dipole moment or density current).
A temporal profile assigned to the vector potential, $\vb A(t)$, does not
produce the same profile in the electric field, $\vb E(t)$, because the latter
contains a time derivative of the former. 
This distinction is important for the RMG perturbation discussed
below. 


For finite systems, the relevant observable is the time-dependent
electronic dipole moment
$\vb d(t)=q_e\int\vb r\,\rho(\vb r,t)\,d^3r$,
from which the polarizability and absorption cross section are obtained
by Fourier analysis of the induced dipole.\cite{yabana-bertsch-1996,jakowski-rmg-tddft-2025}
%----
For periodic systems, however, the position operator, $\vb {\hat r} =\vb r$,
is not compatible with periodic boundary conditions and  the relevant observable is
the induced macroscopic charge-current density
$\vb J(t)$, obtained by averaging the microscopic current density
$\vb j(\vb r,t)$ over the simulation cell.\cite{Luber-2023-position-operator,Resta-polarization-1994,KingSmith-Vanderbilt-polarization-1993}


In summary, it is useful at this point to \myblue{distinguish
the homogeneous (spatially uniform, $\vb q=0$) optical perturbations 
discussed above from probes with finite momentum transfer}.
\myblue{In the homogeneous velocity-gauge representation},
the external field enters through the vector potential in the
kinetic energy via the minimal coupling
$\vb{\hat p}\to\vb{\hat p}-(q_e/c)\vb A$,
\myblue{together with the consistent coupling of nonlocal operators}.
The choice of gauge (length or velocity) determines how
this coupling is represented in the Hamiltonian,
\myblue{while the homogeneous-field approximation determines
the spatial character of the perturbation}.
Typical applications include optical absorption spectra
and laser-driven electron dynamics.
\myblue{Applications using a homogeneous vector potential also include
laser-induced demagnetization when spin--orbit coupling and noncollinear
spin dynamics are included,}\cite{Gross-2015-demagnetization}
\myblue{and time- and angle-resolved photoemission (ARPES).
The latter additionally requires extraction of the outgoing photoelectron
flux to obtain momentum-resolved photoelectron spectra.}\cite{Rubio-angular-2017} 

\myblue{A finite-momentum longitudinal electric probe can instead
be represented by a scalar potential} of the form
\begin{equation}
 \Phi_{\mathrm{ext}}(\vb r,t)
 =\Phi_0 e^{\imath(\vb q\cdot\vb r-\omega t)}
 +\mathrm{c.c.},
 \label{eq:scalar_probe}
\end{equation}
\myblue{with $\vb A=0$. The perturbing potential energy
$q_e\Phi_{\mathrm{ext}}$ is added directly to the Hamiltonian}.
The electric field generated by such a potential,
\myblue{$\vb E=-\nabla\Phi_{\mathrm{ext}}
=-\imath\vb q\Phi_0 e^{\imath(\vb q\cdot\vb r-\omega t)}
+\mathrm{c.c.}$},
is purely longitudinal (parallel to $\vb q$)
and carries no magnetic field.
This type of perturbation probes the density--density
response function $\myblue{\chi_{nn}}(\vb q,\omega)$
and is a natural tool for studying collective excitations
at finite momentum transfer, such as plasmon dispersion
and charge-density oscillations.
\myblue{The subscript distinguishes the density response from
the gauge function $\chi(\vb r,t)$.
Density-response and dielectric-response processing connects
such calculations to EELS loss spectra.}\cite{Weissker-2006,Weissker-2010,JCTC-2020-ebeam,Rubio-2002-eels}

\myblue{When this perturbation is imposed directly in a periodic
supercell with lattice vectors $\vb L_i$, periodicity requires
$\vb q\cdot\vb L_i=2\pi n_i$, where $n_i$ are integers.
For a cell of length $8a$ along one direction, the allowed components
are $q_m=2\pi m/(8a)$, giving a spacing of $1/8$ in units of
the primitive reciprocal vector $2\pi/a$.
Larger supercells permit smaller nonzero momentum transfers.
The $q=0$ scalar plane wave is spatially constant and creates
no electric field; the homogeneous optical limit requires the
separate treatment described above.}

\myblue{A finite-$\vb q$ transverse perturbation can be represented
by Eq.~\eqref{eq:plane_wave_vector_potential} with
$\Phi_{\mathrm{ext}}=0$ and $\vb q\cdot\vb A_0=0$.
For this field,}
\begin{equation}
 \myblue{
 \vb E=\frac{\imath\omega}{c}\vb A_0
 e^{\imath(\vb q\cdot\vb r-\omega t)}+\mathrm{c.c.},
 \qquad
 \vb B=\imath\vb q\times\vb A_0
 e^{\imath(\vb q\cdot\vb r-\omega t)}+\mathrm{c.c.}.}
 \label{eq:em-transverse-fields}
\end{equation}
Although both \myblue{finite-momentum probes} involve the same
spatial factor $e^{\imath\vb q\cdot\vb r}$, they are physically
and computationally distinct:
\myblue{the scalar representation above produces a longitudinal
electric field with no magnetic field, whereas the transverse
vector representation produces a transverse electric field and
an associated magnetic field.
Both admit electromagnetic gauge transformations.
Thus the distinction between the physical probes is not the
presence or absence of gauge freedom.}

Many important spectroscopies and collective excitations involve
finite momentum transfer and lie beyond the dipolar approximation.
These include plasmon dispersion \myblue{as reflected in}
$\varepsilon(\vb q,\omega)$, electron energy-loss spectroscopy
(EELS), where the beam transfers finite $\vb q$ to the sample,
and \myblue{nonresonant} inelastic X-ray scattering (IXS).
\myblue{Nonresonant IXS probes density correlations, so the density
response computed using a finite-$\vb q$ scalar probe is relevant to its
interpretation. Here $\vb q$ is the momentum-transfer wave vector, the
difference between the incident and scattered photon wave vectors;
the scalar probe does not represent the incident X-ray field itself.}\cite{Dixit-2014-IXS}
\myblue{Spatially varying vector potentials instead provide access to
optical response and light--matter coupling beyond the electric-dipole
approximation, as demonstrated in full minimal coupling Maxwell--TDDFT.
That framework treats more general fields than the single transverse
plane-wave component shown in the table.}\cite{Bonafe-Rubio-2025}

\myblue{Magnon spectra and real-time magnon dynamics are obtained using
spin-dependent magnetic perturbations, such as magnetic kicks or
spatially patterned fields that tilt the magnetization, followed by
analysis of the evolving spin magnetization.}\cite{Rubio-2020-magnons,Gross-2020-rt-magnons}
\myblue{These perturbations couple to the spin magnetic moment and
require a dynamical treatment of the spin direction. In this context,
``transverse'' refers to the equilibrium magnetization direction,
whereas in the electromagnetic table it refers to the wave vector
$\vb q$. A finite magnon wave vector does not imply excitation by a
finite-wave-vector electromagnetic plane wave.}
\myblue{At equilibrium, the magnetic dynamic structure factor is related
to the dissipative magnetic susceptibility through the
fluctuation--dissipation relation. Comparison with magnetic INS
additionally requires the appropriate neutron-scattering polarization
factors.}
Table~\ref{tab:gauge_cases} summarizes \myblue{representative electromagnetic
and spin-dependent perturbations} used in TDDFT and their typical applications.
\myblue{The first two rows describe the homogeneous optical
representations relevant here; the finite-$\vb q$ and magnetic/spin-probe
rows provide context for future extensions of the implementation.}

\begin{table}[ht!]
\centering
\footnotesize
\caption{ Representative electromagnetic and spin-dependent
perturbations used in linear-response and real-time TDDFT.
The potentials describe the applied perturbation, with electron
potential energy $q_e\Phi$. In the finite-$\vb q$ rows,
$\theta=\vb q\cdot\vb r-\omega t$; c.c.\ denotes the complex
conjugate. The $B$ column distinguishes the electromagnetic
field $\vb B=\nabla\times\vb A$ from the independently applied
spin field in the last row, whose orbital coupling is omitted.
%-----CUT-----	\myblue{Representative electromagnetic and spin-dependent
%perturbations used in TDDFT}, applicable to both linear-response and
%real-time formulations.
%\myblue{The potentials refer to the applied perturbation;
%$\Phi$ enters the electron Hamiltonian as $q_e\Phi$.
%In the finite-$\vb q$ rows, $\theta =\vb q\cdot\vb r-\omega t$,
%and complex-conjugate terms are included to give real potentials.
%The $B$ column indicates the magnetic field associated with the
%electromagnetic potentials; the auxiliary spin field in the last row
%is identified separately.
%In Coulomb gauge at nonzero $\vb q$, the scalar potential
%represents the longitudinal electric field, while the transverse
%vector potential represents the transverse electric field and
%magnetic field.}
%\myblue{In the magnetic/spin-probe row,
%$\delta\vb B_s$ is an auxiliary field coupled to the spin
%magnetic-moment operator $\hat{\boldsymbol{\mu}}_s$.
%Orbital electromagnetic coupling is omitted in this probe;
%$\delta\vb B_s$ is therefore prescribed independently of
%$\Phi_{\mathrm{ext}}$ and $\vb A_{\mathrm{ext}}$.}
	}
\label{tab:gauge_cases}
\smallskip
\setlength{\tabcolsep}{3pt}
\renewcommand{\arraystretch}{1.15}
\begin{tabular}{@{}
	p{0.145\linewidth}
	p{0.190\linewidth}
	p{0.185\linewidth}
	p{0.120\linewidth}
	p{0.275\linewidth}
	@{}}
\hline
Case & \myblue{Potentials / coupling} & Field character & $B$ field & Typical applications\\
\hline
\raggedright Length gauge, dipole
& \raggedright $\Phi=-\vb r\cdot\vb E(t)$,\newline $\vb A=0$
& \raggedright homogeneous electric field
& no
& \raggedright optical absorption of finite systems, photochemistry
\cite{yabana-bertsch-1996,jakowski-rmg-tddft-2025}\tabularnewline[0.6em]
\raggedright Velocity gauge\myblue{, homogeneous}
& \raggedright $\Phi=0$,\newline $\vb A=\vb A(t)$\newline
\myblue{(uniform in space)}
& \raggedright homogeneous \myblue{electric} field;\newline
\myblue{$\vb E=-(1/c)\dot{\vb A}$}
& no
& \raggedright optical response of solids and periodic systems under PBC
\cite{Bertsch2000,Mattiat-Luber-2022,Hanasaki2023};
\myblue{laser-induced demagnetization (with spin--orbit coupling)\cite{Gross-2015-demagnetization};
	time- and angle-resolved photoemission (ARPES)\cite{Rubio-angular-2017}}\tabularnewline[0.6em]
\raggedright Finite-$\vb q$ scalar probe
& \raggedright \myblue{$\Phi=\Phi_0e^{\imath \theta}+\mathrm{c.c.}$},\newline $\vb A=0$
& \raggedright longitudinal,\newline $\vb E\parallel\vb q$
& no
& \raggedright density response, loss function, EELS, plasmon dispersion,
\myblue{nonresonant IXS}\myblue{\cite{Weissker-2006,Weissker-2010,JCTC-2020-ebeam,Rubio-2002-eels}}\tabularnewline[0.6em]
\raggedright Finite-$\vb q$ \myblue{transverse} EM wave
& \raggedright \myblue{$\Phi=0$},\newline
\myblue{$\vb A=A_0\boldsymbol{\epsilon}e^{\imath \theta } +\mathrm{c.c.,} $}\newline
%\myblue{$\phantom{\vb A={}}+\mathrm{c.c.}$},\newline
$\vb q\cdot\boldsymbol{\epsilon}=0$
& \raggedright transverse EM radiation;\newline
\myblue{$\vb E\perp\vb q$}
& yes
& \raggedright \myblue{optical response and light--matter coupling beyond
the electric-dipole approximation\cite{Bonafe-Rubio-2025}}\tabularnewline
%---
\raggedright \myblue{Magnetic/spin probe}
& \raggedright
\myblue{$\Phi_{\mathrm{ext}}=0$, $\vb A_{\mathrm{ext}}=0$;
\newline
%$\delta\hat H=$\newline
%$-\hat{\boldsymbol{\mu}}_s\cdot\delta\vb B_s(\vb r,t)$} 
	$\delta\hat H=-\hat{\boldsymbol{\mu}}_s\cdot\delta\vb B_s(\vb r,t)$ }
& \raggedright
\myblue{spin-dependent; uniform or spatially varying}
& \raggedright
\myblue{auxiliary spin field}
& \raggedright
\myblue{spin susceptibility, spin precession,
magnon spectra and dynamics; magnetic inelastic neutron
scattering (INS)\cite{Rubio-2020-magnons,Gross-2020-rt-magnons}}
\tabularnewline
%---
\hline
\end{tabular}
\end{table}
